\section{Splitting with No Damping}
\label{Sec:MS-split}
The experiments presented in~\cite{CohenGZ20} and in Section~\ref{Sec:experiments} demonstrate that when standard Min-sum is performed on an SCFG, the algorithm oscillates between low quality solutions and, in many cases, between only two such solutions.
\begin{figure}
\centering
\includegraphics[width=8cm]{two_cycle_examp.pdf}
\vspace{-10pt}
\caption{A symmetric SCFG in which two binary function-nodes were split. The tables show the original costs; each split copy holds half of every entry.}
\vspace{-5pt}
\label{Fig:scfg}
\end{figure}
To investigate this behavior, we examined an example of a lemniscate factor graph, which was presented in~\cite{CohenGZ20}. This factor graph, depicted in Figure~\ref{Fig:scfg}, includes two cycles that were generated from the symmetric split of two function-nodes. As stated in~\cite{CohenGZ20}, the MS algorithm does not converge on this factor-graph. In fact, it alternately visits two solutions, one with cost 300 and the other with cost 1200. However, when damping is used, the algorithm converges to the optimal solution with cost 130.
For the case of a single cycle factor graph,~\citet{ForneyKMT2001} prove that the oscillation between multiple solutions is the result of a minimal inconsistent route that the algorithm follows. Thus, one may think that the oscillation in the case of a lemniscate is also caused by an inconsistent minimal route. In fact, after analyzing the case depicted in Figure~\ref{Fig:scfg}, we arrive at the following observation:
\begin{observation}[Unsynchronized period-2 oscillation]\label{obs:lemniscate}
On the lemniscate SCFG of Figure~\ref{Fig:scfg}, MS converges to a period-2 oscillation between two \emph{solutions}: the optimal solution, whose cost is 130, and a disjoint alternative solution with cost 132. The variable-nodes visit the two solutions out of phase --- in iterations in which $X_1$ and $X_3$ take the value assignments of the optimal solution, $X_2$ takes the value assignment of the alternative solution, and vice versa --- so the decoded joint assignments have costs 300 and 1200.
\end{observation}
\noindent This explains why neither of the two good solutions appears in the reported results: the nodes are not synchronized in their oscillation between them. (A BCT view of the same phenomenon --- the alternating structure of the bottom layers of the middle variable's backtrack cost tree --- is consistent with this observation, but the general account is given in Section~\ref{sec:tworoutes}: once the split function-nodes commit, the decoded assignments follow a synchronous local selection rule whose only limit behaviors are convergence and an alternation between exactly two solutions, and the out-of-phase structure of Observation~\ref{obs:lemniscate} is precisely the bipartite rephasing of Corollary~\ref{cor:bipartite}, under which the two interleaved snapshots of costs 300 and 1200 carry two locally optimal solutions of costs 130 and 132.)
\paragraph{Reading the two branches.}
Observation~\ref{obs:lemniscate} suggests a diagnostic.  If the raw
snapshots of an oscillating MS-SCFG run are two locally optimal solutions
read out of phase, then restricting every variable to the two values it
alternates between --- its \emph{binary menu} --- should contain an
assignment far better than either snapshot, and the way the values are
chosen from the menus should matter.  Section~\ref{Sec:experiments} tests
this: we run MS on the SCFG for 200 iterations, take the assignments
selected at iterations 198 and 199, and choose within the menus either
with MGM~\cite{MaheswaranTBPV04} or exactly, by branch and
bound~\cite{HirayamaY97}.  On every benchmark the chosen assignment is
dramatically better than the raw oscillation, and inverting every
disagreement decision of the MGM choice is significantly worse
(Table~\ref{tab:merge}) --- though still far better than the raw
snapshots, as Corollary~\ref{cor:bipartite} leads one to expect when both
rephasings of the pair are locally optimal.  The reading also has a
limit.  The menus are only as good as the two branches, and on every
benchmark the best choice within them falls short of the split-damped
algorithms of Table~\ref{tab:final_costs}, which resolve the alternation
inside the message passing rather than after it.  We therefore use the
menu selection as evidence for the structure analyzed next, not as a
solver: Section~\ref{sec:tworoutes} explains why the alternation is
between exactly two solutions, and Section~\ref{sec:whydamping} why
damping removes it.
